\documentclass[UTF8,fontset=windows,12pt]{ctexart}
%\setcounter{tocdepth}{1}
%\setcounter{secnumdepth}{4}
%\usepackage{ctex} %若要输入中文请取消注释或者把article改成ctexart
% This first part of the file is called the PREAMBLE. It includes
% customizations and command definitions. The preamble is everything
% between \documentclass and \begin{document}.
%\usepackage[OT1]{fontenc}
%\usepackage{fontspec}
%\setmainfont{Times New Roman}
%\setCJKmainfont{Microsoft YaHei}
\usepackage[left=1.91cm,right=1.91cm,top=2.54cm,bottom=2.54cm]{geometry}
\usepackage{amsfonts}
\usepackage{amsmath}
\usepackage{amssymb}
\usepackage[upint,lcgreekalpha]{stix}
\usepackage[type1]{garamondlibre}
%\usepackage{esint}
\usepackage{pgfplots}
\pgfplotsset{compat=1.18}
\usepackage{tikz}
\usetikzlibrary{arrows}
\usepackage{extarrows}
\usepackage[only,llbracket,rrbracket]{stmaryrd}
\usepackage{titlesec}
\usepackage{bm}
\usepackage{graphicx}
\usepackage{appendix}
\usepackage{tikz-cd}
\usepackage{tikz-3dplot}
\usetikzlibrary{arrows.meta, decorations.pathreplacing, calc, patterns}
\usepackage{float}
\usepackage{mathtools}
\usepackage{amsthm}
\usepackage{verbatim}              % better theorem environments
\usepackage{setspace}
\usepackage{enumitem}
\setenumerate[1]{itemsep=0pt,partopsep=0pt,parsep=\parskip,topsep=5pt}
\setitemize[1]{itemsep=0pt,partopsep=0pt,parsep=\parskip,topsep=5pt}
\setdescription{itemsep=0pt,partopsep=0pt,parsep=\parskip,topsep=5pt}
\usepackage{xcolor}
\usepackage[colorlinks,linkcolor=black,anchorcolor=black,citecolor=black]{hyperref}
%\usepackage{apacite}
\usepackage[numbers]{natbib}
%\usepackage{showkeys}
\usepackage{cases}
\usepackage{fancyhdr}
\usepackage[scaled=0.92]{helvet}	% set Helvetica as the sans-serif font
%\usepackage{upgreek}
%\renewcommand{\rmdefault}{ptm}		% set Times as the default text font
%\usepackage{newtxtext}
\bibliographystyle{plain}
% If AMS-LaTeX is used, it can be loaded before or after mtpro2		
% The following loads metro and defines some common MTPro options [2, 4]
%\usepackage[lite,subscriptcorrection,nofontinfo]{mtpro2}
%\pdfmapfile{+mtpro2.map}

%\usepackage{txfonts}

% various theorems, numbered by section

\makeatletter
\renewenvironment{proof}[1][\proofname]{\par
  \pushQED{\qed}%
  \normalfont \topsep6\p@\@plus6\p@\relax
  \trivlist
  \item[\hskip\labelsep
        \bfseries % 关键修改：斜体改为加粗
        #1\@addpunct{.}]\ignorespaces
}{%
  \popQED\endtrivlist\@endpefalse
}
\makeatother
\newtheoremstyle{kaiti-theorem}   % 样式名称
  {3pt}                           % 上方间距
  {3pt}                           % 下方间距
  {\kaishu\upshape}               % 正文字体：中文楷体 + 英文直立
  {}                              % 缩进量
  {\bfseries\upshape}      % 头部字体（定理名）：中文楷体粗体 + 英文直立粗体
  {.}                             % 定理名后标点
  {0.5em}                         % 定理名后间距
  {}                              % 头部说明（留空）

% 应用自定义样式
\theoremstyle{kaiti-theorem}
\newtheorem{thm}{定理}[section]
\newtheorem{nota}[thm]{记号}
\newtheorem{question}{问题}
\newtheorem*{questionn}{思考}
\newtheorem{exe}{习题}[section]
\newtheorem{assump}[thm]{假设}
\newtheorem*{thmm}{定理}
\newtheorem{defn}{定义}[section]
\newtheorem{lem}[thm]{引理}
\newtheorem*{lemm}{引理}
\newtheorem{prop}[thm]{命题}
\newtheorem*{propp}{命题}
\newtheorem*{claim}{断言}
\newtheorem{cor}[thm]{推论}
\newtheorem{conj}[thm]{猜想}
\newtheorem{rmk}{注记}[section]
\newtheorem{ex}{例}[section]
\newtheorem{pf}{证明}
\newtheorem{problem}{作业题}
\numberwithin{equation}{section}





\DeclareMathOperator{\id}{id}
\DeclareMathOperator*{\esssup}{ess\,sup}
\renewcommand{\Re}{\textbf{Re}\,}  
\renewcommand{\Im}{\textbf{Im}\,}  
\newcommand{\R}{\mathbb{R}}  
\newcommand{\FH}{\mathbb{H}}  
\newcommand{\C}{\mathbb{C}}  
\newcommand{\Z}{\mathbb{Z}}
\newcommand{\X}{\mathbb{X}}
\newcommand{\N}{\mathbb{N}}
\newcommand{\Q}{\mathbb{Q}}
\newcommand{\T}{\mathbb{T}}
\newcommand{\TT}{\mathcal{T}}
\newcommand{\TP}{\overline{\partial}{}}
\newcommand{\TJ}{\langle\TP\rangle}
\newcommand{\TL}{\overline{\Delta}{}}
\newcommand{\curl}{\text{curl }}
\newcommand{\dive}{\text{div }}
\newcommand{\bd}[1]{\mathbf{#1}}  % for bolding symbols
\newcommand{\bds}[1]{\boldsymbol{#1}}  % for bolding symbols
\newcommand{\RR}{\mathcal{R}}      % for Real numbers
\newcommand{\sh}{\mathcal{S}}  
\newcommand{\sph}{\mathbb{S}}  
\newcommand{\Om}{\Omega}
\newcommand{\OmT}{{\Omega_T}}
\newcommand{\ut}{{U_T}}
\newcommand{\vp}{\varphi}
\newcommand{\Th}{\Theta}
\newcommand{\Lam}{\Lambda}
\newcommand{\Omb}{{\overline{\Omega}}}
\newcommand{\OmbT}{{\overline{\Omega_T}}}
\newcommand{\ub}{{\overline{U}}}
\newcommand{\ubt}{{\overline{U_T}}}
\newcommand{\q}{\quad}
\newcommand{\p}{\partial}
\newcommand{\pb}{\boldsymbol{\partial}}
\newcommand{\lap}{\Delta}
\newcommand{\dd}{\mathfrak{D}}
\newcommand{\DD}{\mathcal{D}}
\newcommand{\den}{{\bar\rho}_{0}}
\newcommand{\Dt}{{\p}_{t}+v^{k}{\p}_{k}}
\newcommand{\col}[1]{\left[\begin{matrix} #1 \end{matrix} \right]}
\newcommand{\noi}{\noindent}
\newcommand{\nab}{\nabla}
\newcommand{\cnab}{\overline{\nab}}
\newcommand{\cp}{\overline{\partial}{}}
\newcommand{\dx}{\,\mathrm{d}x}
\newcommand{\dxx}{\,\mathrm{d}\bds{x}}
\newcommand{\xxi}{\boldsymbol{\xi}}
\newcommand{\eeta}{\boldsymbol{\eta}}
\newcommand{\dxxi}{\,\mathrm{d}\boldsymbol{\xi}}
\newcommand{\deeta}{\,\mathrm{d}\boldsymbol{\eta}}
\newcommand{\dxi}{\,\mathrm{d}\xi}
\newcommand{\deta}{\,\mathrm{d}\eta}
\newcommand{\dy}{\,\mathrm{d}y}
\newcommand{\dyy}{\,\mathrm{d}\bds{y}}
\newcommand{\dz}{\,\mathrm{d}z}
\newcommand{\dzz}{\,\mathrm{d}\bds{z}}
\newcommand{\dph}{\,\mathrm{d}\varphi}
\newcommand{\dt}{\,\mathrm{d}t}
\newcommand{\dr}{\,\mathrm{d}r}
\newcommand{\dtt}{\,\mathrm{d}\tau}
\newcommand{\dS}{\,\mathrm{d}S}
\newcommand{\dss}{\,\mathrm{d}\sigma}
\newcommand{\dmu}{\,\mathrm{d}\mu}
\newcommand{\dnu}{\,\mathrm{d}\nu}
\newcommand{\dV}{\,\mathrm{d}V}
\newcommand{\ds}{\,\mathrm{d}s}
\newcommand{\drr}{\,\mathrm{d}\rho}
\newcommand{\dist}{\text{dist }}
\newcommand{\vol}{\text{vol}\,}
\newcommand{\volo}{\text{vol}(\Omega)}
\newcommand{\spt}{\text{Spt}\,}
\newcommand{\Tr}{\text{Tr}\,}
\newcommand{\lee}{\langle}
\newcommand{\ree}{\rangle}
\newcommand{\xee}{\langle\xi\rangle}
\newcommand{\xxee}{\langle\boldsymbol{\xi}\rangle}
\newcommand{\xeta}{\langle\eta\rangle}
\newcommand{\xeeta}{\langle\boldsymbol{\eta}\rangle}
\newcommand{\kk}{\kappa}
\newcommand{\lam}{\lambda}
\newcommand{\io}{\int_{\Omega}}
\newcommand{\iu}{\int_{U}}
\newcommand{\is}{\int_{\Sigma}}
\newcommand{\ipo}{\int_{\partial\Omega}}
\newcommand{\ipu}{\int_{\partial U}}
\newcommand{\ig}{\int_{\Gamma}}
\newcommand{\ir}{\int_{\R}}
\newcommand{\ird}{\int_{\R^d}}
\newcommand{\irn}{\int_{\R^n}}
\newcommand{\fpi}{\frac{1}{(2\pi)^{\frac{d}{2}}}}
\newcommand{\ddt}{\frac{\mathrm{d}}{\mathrm{d}t}}
\newcommand{\ddx}{\frac{\mathrm{d}}{\mathrm{d}x}}
\newcommand{\eps}{\varepsilon}
\providecommand{\jump}[1]{\left\llbracket #1 \right\rrbracket }
\providecommand{\len}[1]{\lee #1 \ree }
\providecommand{\ino}[1]{\left\| #1 \right\| }
\providecommand{\bno}[1]{\left| #1 \right| }
\newcommand{\red}{\textcolor{red}}
\newcommand{\green}{\textcolor{green}}
\newcommand{\blue}{\textcolor{blue}}
\newcommand{\nnr}{{[n+1]}}
\newcommand{\nnn}{{[n]}}
\newcommand{\nnl}{{[n-1]}}
\newcommand{\nnll}{{[n-2]}}	
\newcommand{\mmr}{{[m+1]}}
\newcommand{\mmm}{{[m]}}
\newcommand{\mml}{{[m-1]}}
\newcommand{\mmll}{{[m-2]}}

%---------------Special variables--------------
\newcommand{\CC}{\mathfrak{C}}  
\newcommand{\NN}{\mathbf{N}}
\newcommand{\LH}{\mathcal{L}}
\newcommand{\HH}{\mathcal{H}}  
\newcommand{\EE}{\mathcal{E}}  
\newcommand{\FT}{\mathcal{F}}  
\newcommand{\fT}{\mathbf{T}}  
\newcommand{\FA}{\mathbf{A}}  
\newcommand{\MH}{\mathcal{M}}
\newcommand{\NH}{\mathcal{N}}
\newcommand{\PH}{\mathcal{P}}
\newcommand{\VH}{\mathcal{V}}
\newcommand{\XX}{\mathfrak{X}}
\newcommand{\xx}{{\bds{x}}}
\newcommand{\bx}{\mathbf{x}}
\newcommand{\by}{\mathbf{y}}
\newcommand{\bp}{\mathbf{p}}
\newcommand{\bq}{\mathbf{q}}
\newcommand{\pp}{{\bds{p}}}
\newcommand{\qq}{{\bds{q}}}
\newcommand{\uu}{\mathbf{u}}
\newcommand{\ww}{\mathbf{w}}
\newcommand{\vv}{\mathbf{v}}
\newcommand{\yy}{{\bds{y}}}
\newcommand{\zz}{{\bds{z}}}
\newcommand{\zo}{\mathbf{0}}
\newcommand{\ee}{\mathbf{e}}
\newcommand{\fa}{\mathbf{a}}
\newcommand{\fb}{\mathbf{b}}
\newcommand{\fc}{\mathbf{c}}
\newcommand{\ff}{\bds{f}}
\newcommand{\fr}{\mathbf{r}}
\newcommand{\fh}{\mathbf{h}}
\newcommand{\fm}{\mathbf{m}}
\newcommand{\fg}{\mathbf{g}}
\newcommand{\FF}{\mathbf{F}}
\newcommand{\GG}{\mathbf{G}}
\newcommand{\BB}{\mathbf{B}}



\begin{document}
\title{\LARGE{\bf  2026年秋季学期~常微分方程~作业三}}
\author{线性常微分方程组}
\date{截止时间：2026.10.12下课前（电子版截至当天23:59:59）}

\maketitle

\textbf{禁止抄袭或使用AI直接生成答案PDF，被抓到者此次作业得零分。}

%考虑如下形式的一阶常微分方程的初值问题
%\begin{equation}\label{2-0: IVP}\tag{2.0.1}
% \xx'(t)=\ff(t, \xx(t)), \qquad \xx(t_0)=\xx_0, 
%\end{equation}
%其中 \(I\subset\R\) 是含有 \(t_0\) 的区间, 未知函数\(\xx:I\to\R^d\), 给定函数 \(\ff:D\to\R^d\) 的定义域\(D\) 是 \(\R\times\R^d\) 中的开集, 并且 \((t_0, \xx_0)\in D\). 

\begin{problem}[习题3.1.1]
求解方程组 $\xx'(t)=A\xx(t)$, 其中系数矩阵 $A$ 分别如下。
\begin{align*}
(1)~A=\begin{pmatrix} 0&0&1\\0&1&0\\1&0&0 \end{pmatrix}\qquad\qquad (4)~A=\begin{pmatrix}1&9&1\\0&9&-8\\0&1&0 \end{pmatrix}\qquad\qquad (5)~A=\begin{pmatrix} 2&-1&-1\\2&-1&-2\\-1&1&2 \end{pmatrix}
\end{align*}

注：本题系数矩阵的特征值都是整数，如果算出奇怪数字肯定是算错了。
\end{problem}

\begin{problem}[习题3.1.2，选做]
考虑变系数齐次线性方程组的初值问题 $\xx'(t)=A(t)\xx(t)$, $\xx(0)=\xx_0$. 其中 $A(t)=[a^{ij}(t)]_{d\times d}$ 各个元素都是\underline{有界}连续函数。变量 $t\geqslant 0$，即我们不考虑时间倒向演化。
\begin{itemize}
\item [(1)] 若全体$a^{ij}$皆为非负函数，且初值$\xx^0$的各分量也非负，证明解的各个分量保持非负。
\item [(2)] 若把(1)的条件改为对全体$i\neq j$有$a^{ij}$皆为非负函数，其它条件不变，则(1)的结论是否成立？
\end{itemize}

提示：(1) 可以由 (2) 得到，若寻求直接证明，则可以考虑其 Picard 迭代序列满足的积分方程；(2) 利用 $a_{ii}(t)$ 的有界性，去充分大的常数 $C>0$ 使得 $A(t)+CI_d$ 的所有元素皆为非负。
\end{problem}

\begin{problem}[习题3.1.4] 
对复方阵 $A$, 定义 $s(A)$ 是 $A$ 全体特征值实部的最大值。
\begin{itemize}
\item [(1)] 证明：对任意 $\eps>0$，存在与 $\eps$ 有关的常数 $C_\eps>0$ 使得对任意 $t\geqslant 0$ 成立 $\|e^{tA}\|\leqslant C_\eps e^{(s(A)+\eps)t}$.
\item [(2)] 证明：存在不依赖 $\eps$ 的常数 $C>0$ 使得对任意 $t\geqslant 0$ 成立 $\|e^{tA}\|\leqslant C e^{s(A)t}$，当且仅当全体满足 $\Re \lam=s(A)$ 特征值所对应的 Jordan 块均为一阶。
\item [(3)] 证明：$\sup\limits_{t\geqslant 0}\|e^{tA}\|<+\infty$ 当且仅当 $A$ 的全体特征值实部皆为非正，且全体实部为零的特征值对应的 Jordan 块皆为一阶。
\end{itemize}
注：这里的矩阵范数定义为 $\|M\|=\sup\limits_{|\xx|=1}|M\xx|$，它满足 $|M\xx|\leqslant \|M\||\xx|$, $\|MN\|\leqslant \|M\|\|N\|$ 以及 $\|M+N\|\leqslant \|M\|+\|N\|$.  本题应先对单个 Jordan 块讨论，并用到命题3.1.1的结论。
\end{problem}


\begin{problem}[习题3.2.1]
解方程组 $\xx'(t)=A\xx(t) + \ff(t)$, 其中系数矩阵 $A$, 源项 $\ff$, 初值 $\xx(0)$ 如下。
{\small\begin{align*}
(1)&~A=\begin{pmatrix} 2&-1&-1\\2&-1&-2\\-1&1&2  \end{pmatrix},~~\ff=\begin{pmatrix}2t\\0\\0\end{pmatrix},~~\xx(0)=\begin{pmatrix}1\\0\\0\end{pmatrix}.\qquad (2)~A=\begin{pmatrix} -1&-1&0\\0&-1&-1\\0&0&-1 \end{pmatrix},~~\ff=\begin{pmatrix}t^2\\2t\\t\end{pmatrix},~~\xx(0)=\begin{pmatrix}0\\0\\0\end{pmatrix}.
\end{align*}}
\end{problem}


\begin{problem}[习题3.2.2]
设常系数方阵 $A$ 的全体特征值的实部皆为负，设向量值连续函数 $\fg$ 有界。令 $$\xx_b(t):=\int_{-\infty}^t e^{(t-s)A} \fg(s)\ds.$$ 证明如下结论。
\begin{itemize}
\item [(1)] 定义 $\xx_b$ 的积分绝对收敛，且是非齐次线性方程组 $\xx'(t)=A\xx(t)+\fg(t)$ 的一个有界整体解。（提示：用习题 3.1.4.）
\item [(2)] (1) 中所述的方程组在$\R$上的有界整体解是唯一的。
\item [(3)] 若 $\fg$ 是以 $T>0$ 为周期的函数，则 $\xx_b$ 也是。
\item [(4)]  (1) 中所述的方程组的其他任意解 $\xx(t)$ 都满足：当 $t\to+\infty$ 时，成立 $|\xx(t)-\xx_b(t)|\to 0$.
\end{itemize}
\end{problem}

\begin{problem}[习题3.2.3]
设 $A(t):[0,+\infty)\to \R^{d\times d}$ 是连续系数方阵，且 $\int_0^{+\infty}\|A(t)\|\dt<+\infty$. 设 $\Phi$ 是满足 $\Phi'(t)=A(t)\Phi(t)$, $\Phi(0)=I_d$ 的标准基解矩阵。证明 $\Phi(t)$ 在 $[0,+\infty)$ 一致有界，且存在可逆方阵 $\Phi_\infty$ 使得当 $t\to+\infty$ 时有 $\|\Phi(t)-\Phi_\infty\|\to 0$. 进而该齐次方程组的每个解都有极限，且从初值映射到极限值的映射是 $\R^d$ 上的线性同构。

提示：把微分方程转化为积分方程即可看出为什么加了条件 $\int_0^{+\infty}\|A(t)\|\dt<+\infty$，然后思考极限方阵 $\Phi_\infty$ 的可逆性是怎么来的。
\end{problem}


\end{document}
